% !TeX program = xelatex
\documentclass[12pt,a4paper]{article}
\usepackage[margin=1in]{geometry}
\usepackage[UTF8,fontset=fandol]{ctex}
\usepackage{amsmath,amssymb,booktabs,array,graphicx,subcaption,enumitem,caption,titlesec,color,bm,multirow}
\usepackage[dvipsnames]{xcolor}
\usepackage{placeins,xpatch,float,listings,makecell,longtable,hyperref,indentfirst}
\usepackage[ruled,linesnumbered]{algorithm2e}
\usepackage{tikz}
\usetikzlibrary{shapes.geometric,arrows.meta,positioning,calc}
\usepackage{tocloft}
\SetAlgoNlRelativeSize{-1}
\SetNlSty{}{}{}
\captionsetup{font={normalsize,rm,bf},labelfont={rm,bf},textfont={rm,bf},labelsep=quad}
\captionsetup[table]{position=top}
\captionsetup[figure]{position=bottom,skip=6pt}
\hypersetup{colorlinks=true,linkcolor=black,citecolor=black,filecolor=black,urlcolor=black}
\setmainfont{Times New Roman}
\definecolor{background}{rgb}{0.95,0.95,0.92}
\definecolor{keyword}{rgb}{0,0,0}
\definecolor{comment}{rgb}{0.0,0.5,0.0}
\definecolor{string}{rgb}{0.6,0.1,0.7}
\definecolor{line_numbers}{rgb}{0.25,0.25,0.25}
\lstset{language=Matlab,backgroundcolor=\color{white},basicstyle=\ttfamily\footnotesize,breakatwhitespace=false,breaklines=true,commentstyle=\color{comment}\itshape,extendedchars=true,frame=single,keepspaces=true,keywordstyle=\color{blue}\bfseries,numbers=left,numbersep=10pt,numberstyle=\color{line_numbers},rulecolor=\color{black},showspaces=false,showstringspaces=false,showtabs=false,stringstyle=\color{string},tabsize=4}
\newcommand{\jiacu}[1]{\textbf{\bfseries\heiti #1}}
\newcommand{\upcite}[1]{\cite{#1}}
\newcommand{\appendixtocsub}[1]{\addcontentsline{toc}{subsection}{\protect\numberline{}#1}}
\newcommand{\appentry}[2]{附录\makebox[1.5em][l]{#1}#2}

%================== 图片路径 ==================
\graphicspath{{figures/}}

%================== 章节编号 ==================
\renewcommand\thesection{\chinese{section}}
\renewcommand\thesubsection{\arabic{section}.\arabic{subsection}}
\renewcommand\thesubsubsection{\arabic{section}.\arabic{subsection}.\arabic{subsubsection}}
\renewcommand{\thetable}{\arabic{section}-\arabic{table}}
\renewcommand{\thefigure}{\arabic{section}-\arabic{figure}}
\renewcommand{\theequation}{\arabic{section}-\arabic{equation}}
\counterwithin*{figure}{section}
\counterwithin*{table}{section}
\counterwithin*{equation}{section}

%================== 章节标题格式 ==================
\titleformat{\section}[block]
  {\centering\zihao{4}\bfseries\heiti}
  {第\thesection 章}{1em}{}
\titlespacing{\section}{0pt}{18pt}{18pt}
\setlength{\parindent}{2em}

\titleformat{\subsection}
  {\zihao{-4}\bfseries\heiti}
  {\thesubsection}{1em}{}

\titleformat{\subsubsection}
  {\zihao{-4}\bfseries\heiti}
  {\thesubsubsection}{1em}{}

%================== 目录格式 ==================
\renewcommand{\contentsname}{目\quad 录}
\renewcommand{\cfttoctitlefont}{\hfil\fontsize{18pt}{24pt}\bfseries\heiti\selectfont}
\renewcommand{\cftaftertoctitle}{\hfil}
\setlength{\cftbeforetoctitleskip}{0pt}
\setlength{\cftaftertoctitleskip}{20pt}

\renewcommand{\cftsecfont}{\normalsize\heiti}
\renewcommand{\cftsecpagefont}{\normalsize\mdseries}
\renewcommand{\cftsecpresnum}{第}
\renewcommand{\cftsecaftersnum}{章}
\renewcommand{\cftsecaftersnumb}{\quad}
\setlength{\cftsecnumwidth}{3.2em}
\renewcommand{\cftsecleader}{\normalfont\cftdotfill{\cftdotsep}}
\setlength{\cftbeforesecskip}{6pt}

\renewcommand{\cftsubsecfont}{\small}
\renewcommand{\cftsubsecpagefont}{\small}
\setlength{\cftsubsecindent}{3em}
\setlength{\cftsubsecnumwidth}{2.2em}
\renewcommand{\cftsubsecleader}{\normalfont\cftdotfill{\cftdotsep}}
\setlength{\cftbeforesubsecskip}{3pt}

\renewcommand{\cftsubsubsecfont}{\small}
\renewcommand{\cftsubsubsecpagefont}{\small}
\setlength{\cftsubsubsecindent}{4em}
\setlength{\cftsubsubsecnumwidth}{2.2em}
\renewcommand{\cftsubsubsecleader}{\normalfont\cftdotfill{\cftdotsep}}
\setlength{\cftbeforesubsubsecskip}{2pt}

\renewcommand{\cftdotsep}{1.5}
\setcounter{tocdepth}{3}
\setcounter{secnumdepth}{3}
\allowdisplaybreaks

%================== 流程图节点样式 ==================
\tikzset{
  flowbox/.style={rectangle,rounded corners,draw=black,thick,align=center,
    text width=10.5cm,minimum height=1.0cm,font=\small,fill=white},
  flowarrow/.style={-{Stealth[length=3mm]},thick}
}

\begin{document}
\zihao{-4}\songti\linespread{1}\selectfont
\hypersetup{pageanchor=false}

%================== 封面 ==================
\begin{titlepage}
\thispagestyle{empty}
\centering
\vspace*{8em}
{\zihao{3}\heiti 工件 X 光透射检测数据的可分结构分析\\[0.5em]与正则化判别建模\par}
\vspace{6em}
{\zihao{-4}\songti 作者：葛旭东、李悦昊、黎子豪\par}
\vspace{2em}
{\zihao{-4}\songti 院系：新闻与传播学院\par}
\vfill
\end{titlepage}

%================== 摘要 ==================
\thispagestyle{empty}
\begin{center}
	{\zihao{4}\heiti 摘\quad 要}
\end{center}

X 光透射检测能在不破坏工件的前提下反映其内部材质与结构特征，是工业产线上常用的质量检测手段。当同一产线上流转多种型号的工件时，能否仅凭检测数据自动判别型号，直接关系到分拣效率与质量追溯的可靠性。本文针对某工厂 A, B, C 三类工件的分类需求，构建了一套以\jiacu{正则化二次判别分析}为主体、多个统计判别模型相互复核的可解释分类方法。每件工件在 155 个固定检测位置上给出一组连续测量值，样本共 7978 件，其中 A 类 3572 件, B 类 2643 件, C 类 1763 件，另有 D1 至 D10 共 10 件未知工件待判别。

特征分析表明，三类工件的差异由两部分构成。整体透射水平上，三类均值依次为 A 类 7.0291, B 类 6.2109, C 类 3.8578，C 类与 A, B 两类之间存在明显的水平间隔。曲线波动程度上，A 类的样本内标准差约为 B 类的两倍，我们由此断定出这一方差结构差异恰是区分A与B的关键。随后的主成分分析进一步显示第一主成分独占 54.17\% 的方差，三类沿该综合强度方向依次排开。

在分类建模阶段，本文按假设强度递增的顺序建立最近质心, 高斯朴素贝叶斯, 正则化 LDA, 正则化 QDA 与 PCA-QDA 五类模型，并在分层五折交叉验证下用总准确率, 平衡准确率, 宏 F1 与混淆矩阵加以比较。为避免验证集信息泄漏，标准化只在每个训练折内部估计。结果与特征分析吻合：越能利用方差信息的模型预测越准。\jiacu{正则化 QDA 的总准确率、平衡准确率与宏 F1 均为 1.0000}，高斯朴素贝叶斯为 0.9917，其余模型在 0.95 附近。C 类召回率在所有模型中均为 1.0000，误判几乎全部集中在 A 与 B 之间。本文以 QDA 作主分类器，其余四个模型用于一致性复核。

对 10 件未知工件，五个模型分别判别后取众数作为结论，并以模型一致数衡量可信度。判别结果为 \jiacu{D1、D2、D3、D10 属于 A 类，D4、D5、D6 属于 B 类，D7、D8、D9 属于 C 类}，五个模型对每件工件均给出一致判别。机制分析进一步表明，A 与 B 在样本内标准差上几乎不重叠，以阈值 $s^*=1.105855$ 即可将二者完全分开，其中 A 类最小值 1.139050 高于 B 类最大值 1.072661。先验, 收缩参数与降维维数三组敏感性分析均表明该结论稳定。考虑到真实产线可能出现设备漂移或新型号，本文还给出基于模型一致数与 QDA 后验概率的拒判规则，以及基于 Hotelling $T^2$ 的批次漂移监控方案。

\bigskip

\noindent
\jiacu{关键词：}工件分类；X 光检测；正则化二次判别分析；主成分分析；交叉验证
\clearpage
\hypersetup{pageanchor=true}
\pagenumbering{arabic}
\pagestyle{plain}

%================== 正文 ==================
\section{问题重述}
\subsection{问题背景}

X 光透射检测是工业制造中常用的无损检测手段。射线穿过工件时强度按指数规律衰减，衰减量取决于材料的密度, 厚度与成分，因此透射后测得的强度携带了工件内部结构的信息\upcite{mery2015,zhanghui2022}。在一个工件上沿固定路线布置 155 个检测位置，将各位置的测量值按序排列，即可得一条反映该工件透射特性的离散曲线。同一型号工件由相同工艺生产，检测曲线形态相近，不同型号在材料或结构上的差异则让曲线在整体水平或局部波动上呈现系统性的不同。

工厂生产线上常常同时流转多种型号的工件。若能依据检测曲线自动判别型号，即可支撑分拣, 质量追溯与产线监控，减少人工核对的成本与差错。本文的研究对象是某工厂的三类工件 A, B, C，每件工件给出 155 个检测位置上的连续测量值。从数据初步观察来看，三者的检测曲线既有明显可分的部分，亦有容易混淆之处，需用合适的统计模型刻画其可分结构，进而对未知工件给出可靠判别。


\subsection{问题提出}

由题易知 A, B, C 三类工件的检测数据与 10 件未知类型工件的检测数据。查表易知表第一列为工件编号，其余 155 列为该工件在各检测位置上的测量值。据题意，需完成三项工作，建立模型分析三类工件检测数据的特征，由此建模型判别 10 件未知工件的类别，最后对模型作出评价。


\section{问题分析}

本题的核心是一个高维多分类问题。每件工件由 155 维连续向量描述，而类别仅有三种。维度远高于类别数，意味着原始空间中可能存在大量冗余，真正用于区分类别的方向也许只有少数几个。循着这一思路，本文拟先揭示三类工件在 155 维空间中依靠什么结构彼此分开，再据此选取与数据特征匹配的判别模型，最后检验判别结论是否稳健。

本文拟先进行特征分析，即三类工件的差异主要体现在整体透射水平上，还是体现在曲线沿检测位置的波动形态上？若是前者，需刻画均值或类中心的模型，若为后者，模型则必须具备捕捉方差与协方差结构的能力。第二个问题关乎差异的分布。即类别间的差异集中于个别检测位置，还是沿整条曲线呈系统性平移？若为后者，155 维信息可以投影到少数综合方向而不丧失判别信息，降维在提供可视化的同时亦能改善高维协方差估计的稳定性。围绕这两个问题，本文从三类工件的全局统计特征出发，以逐位置单因素方差分析定位区分力最强的检测位置，再以主成分分析揭示数据的低维流形结构。

模型对数据分布的先验假设应当与特征分析的实证发现相匹配，同时判别过程须保持可解释与可复核。按照这一标准，本文选取假设强度递增的五个统计判别模型构成对比谱系。最近质心只比较到各类中心的距离，朴素贝叶斯引入各检测位置的边缘分布及其方差，线性判别分析假设三类协方差相同从而产生线性判别边界，二次判别分析放松该假设以允许各类协方差异质，PCA 降维后的二次判别则用于检验高维空间中的结论在压缩表示下是否依然成立。将五个模型置于统一的交叉验证框架下比较，既可遴选最契合数据结构的模型，亦可借助模型间预测的一致或分歧来量化单次判别的置信程度。鉴于三类样本量不均衡，评价指标在总准确率之外需引入对各类等权的平衡准确率，避免样本量最大的 A 类主导整体评估。

在稳健性检验上，本文将于文末讨论所用模型的判别力源于何种统计结构，并检验该结构在不同超参数设定与验证协议下是否稳定，以及当实际生产环境出现设备漂移或新工件型号时，模型能否给出可靠的拒判信号而非静默误判。为此，本文在分类性能报告之外补充机制诊断与敏感性分析，并给出面向产线部署的拒判规则与批次漂移监控方案。上述从特征分析到分类判别再到模型评价的整体研究路径如图~\ref{fig:flow} 所示。

\begin{figure}[H]
\centering
\includegraphics[height=0.74\textheight]{08_modeling_flowchart}
\caption{本文建模流程图}
\label{fig:flow}
\end{figure}

\section{模型假设与符号说明}
\subsection{模型假设}

为使建模过程清晰可控，本文作如下假设。

\begin{enumerate}[label=(\arabic*),leftmargin=2.4em]
\item 同一类工件由相同工艺生产，其 155 维检测向量服从同一分布，类内差异由材质波动与测量噪声引起。
\item 训练样本的类别标签准确，A, B, C 三类划分无误。
\item 各检测位置的测量在同一标定条件下获得，训练数据与未知数据来自同一检测系统，量纲一致。
\item 训练集中三类样本的数量比例近似反映生产线的真实比例，先验概率可由样本比例估计。
\item 未知工件属于 A, B, C 三类之一。该假设用于本题判别.在模型推广部分另行讨论出现新类型时的拒判处理。
\item 训练样本可近似视为独立同分布，交叉验证的随机划分不会将强相关的同批次样本同时分入训练与验证，保证性能估计有效。
\end{enumerate}

\subsection{符号说明}

\begin{longtable}{cl}
    \toprule
    符号 & 意义 \\
    \midrule
    \endfirsthead
    \multicolumn{2}{c}{{\tablename\ \thetable{} 续页}} \\
    \toprule
    符号 & 意义 \\
    \midrule
    \endhead
    \midrule
    \multicolumn{2}{r}{{续下页}} \\
    \endfoot
    \bottomrule
    \endlastfoot
    $N$ & 训练样本总数，$N=7978$ \\
    $p$ & 检测位置数，$p=155$ \\
    $G$ & 类别数，$G=3$ \\
    $\mathcal{G}$ & 类别集合 $\{A,B,C\}$ \\
    $n_g$ & 类别 $g$ 的样本数 \\
    $\mathbf{x}_i$ & 第 $i$ 件工件的原始检测向量，$\mathbf{x}_i\in\mathbb{R}^{155}$ \\
    $x_{ij}$ & 第 $i$ 件工件在第 $j$ 个检测位置的测量值 \\
    $y_i$ & 第 $i$ 件工件的类别标签 \\
    $\mu_j,\,s_j$ & 第 $j$ 个检测位置在训练集上的均值与标准差 \\
    $z_{ij},\,\mathbf{z}_i$ & 标准化后的测量值与样本向量 \\
    $\bar{x}_i,\,s_i,\,R_i,\,IQR_i$ & 单件工件检测曲线的均值, 样本内标准差, 极差与四分位距 \\
    $\boldsymbol{\mu}_g$ & 类别 $g$ 在标准化空间的类中心 \\
    $\boldsymbol{\Sigma},\,\boldsymbol{\Sigma}_g$ & 共享协方差矩阵与类别 $g$ 的协方差矩阵 \\
    $\pi_g$ & 类别 $g$ 的先验概率 \\
    $\lambda,\,\alpha$ & LDA 与 QDA 的协方差收缩参数 \\
    $\delta_g^{(\cdot)}$ & 各模型对类别 $g$ 的判别函数 \\
    $c_{gh}$ & 混淆矩阵元素，真实类别 $g$ 被判为 $h$ 的样本数 \\
    $ACC,\,BA$ & 总准确率与平衡准确率 \\
    $F_j,\,\eta_j^2$ & 第 $j$ 个检测位置的方差分析 F 统计量与效应量 \\
    $\lambda_k,\,\mathbf{v}_k,\,R_k$ & 第 $k$ 个主成分的特征值, 载荷向量与解释率 \\
    $s^{*}$ & 区分 A, B 两类的样本内标准差最优阈值 \\
    $T^2$ & 用于批次漂移监控的 Hotelling 统计量 \\
\end{longtable}

\section{数据预处理与特征分析}
\subsection{数据读取与质量控制}

四表分别对应 A, B, C 三类与未知类型工件。读取时将第一列作为工件编号，其余 155 列作为数值特征。合并三类已知工件后得到训练矩阵 $\mathbf{X}\in\mathbb{R}^{7978\times155}$ 与标签向量 $\mathbf{y}$，其中 A, B, C 三类分别为 3572, 2643, 1763 件。逐类逐元素统计缺失值，三类的缺失值数量均为 0，未知数据的列数与训练数据一致，无需插补。

\subsection{标准化处理}

不同检测位置的量值水平与波动程度各异，若直接使用原始值，量值较大的位置会主导距离与协方差计算。为消除量纲影响，对每个检测位置做Z分数标准化。第 $j$ 个位置在训练集上的均值与标准差为
\begin{equation}
\mu_j=\frac{1}{N}\sum_{i=1}^{N}x_{ij},\qquad
s_j=\sqrt{\frac{1}{N}\sum_{i=1}^{N}(x_{ij}-\mu_j)^2},
\end{equation}
标准化后的测量值为
\begin{equation}
z_{ij}=\frac{x_{ij}-\mu_j}{s_j}.
\end{equation}
在交叉验证中，$\mu_j$ 与 $s_j$ 只能由当前训练集估计后再作用于验证集。

\subsection{三类工件的全局统计特征}

将单件工件的 155 个检测值视为一条曲线，对第 $i$ 件工件定义四个曲线统计量。均值 $\bar{x}_i=\frac{1}{p}\sum_j x_{ij}$ 反映整体透射水平，样本内标准差 $s_i$ 反映曲线沿检测位置的波动幅度，极差 $R_i$ 与四分位距 $IQR_i$ 分别从最极端两点与中间一半数据的角度补充刻画离散程度。前一个统计量回答工件整体偏亮还是偏暗，后三个统计量回答曲线是平缓还是起伏。由此信息可判断三类工件间的差异落在水平上还是波动上。三类工件的总体统计特征列于表~\ref{tab:profile}。

\begin{table}[H]
\centering
\caption{三类工件的全局统计特征}
\label{tab:profile}
\resizebox{\textwidth}{!}{%
\begin{tabular}{cccccccc}
\toprule
类别 & 样本数 & 全体均值 & 全体标准差 & 样本均值标准差 & 样本内标准差均值 & 最小值 & 最大值 \\
\midrule
A & 3572 & 7.0291 & 1.4956 & 0.2040 & 1.4831 & $-0.0029$ & 14.4913 \\
B & 2643 & 6.2109 & 0.8509 & 0.3236 & 0.7847 & 2.6257 & 10.3303 \\
C & 1763 & 3.8578 & 0.8300 & 0.2523 & 0.7868 & 0.6268 & 8.0281 \\
\bottomrule
\end{tabular}}
\end{table}

查表易知，从水平维度看，三类整体均值呈 A 高于 B 高于 C 的次序，A 与 B 相差 0.8182，B 与 C 相差 2.3531。C 与另外两类之间的水平间隔约为 A 与 B 之间间隔的三倍，单凭整体水平已足以将 C 辨识出来，而 A 与 B 的水平过于接近，仅靠均值难以将二者区分。从波动维度看，A 类的样本内标准差均值为 1.4831，约是 B 类 0.7847 与 C 类 0.7868 的两倍，这表明 A 类工件的检测曲线起伏明显比 B 和 C 剧烈。那么我们不难清楚，即便 A 与 B 的整体水平相近，二者在曲线波动幅度上仍存在系统性差异，只要模型能够利用方差信息，就有可能将 A 与 B 分开。图~\ref{fig:mean} 给出三类工件的平均检测曲线及其 $\pm1$ 标准差范围。

\begin{figure}[H]
\centering
\includegraphics[width=1.0\textwidth]{01_mean_curves}
\caption{三类工件平均检测曲线及 $\pm1$ 标准差阴影}
\label{fig:mean}
\end{figure}

从图~\ref{fig:mean} 可以看出，C 类曲线整体位于下方且与 A, B 两类几乎不重叠，A 类曲线位于上方且阴影最宽，B 类介于两者之间且波动较窄。这与表~\ref{tab:profile} 的统计量一致。C 类容易识别，而 A 与 B 的区分需要同时利用水平差异与波动差异。

\subsection{关键检测位置的方差分析}

全局统计特征说明三类整体可分，但尚未回答差异是集中于少数检测位置还是均匀分布在整条曲线上。因此本文接下来使用单因素方差分析做探索，即将同一检测位置上三类的取值看作三组样本，比较组间差异与组内波动的相对大小。若组间差异远超组内波动，说明该位置的取值主要由类别决定，区分力强。

对每个位置 $j$ 分别做单因素方差分析。设类别 $g$ 在位置 $j$ 上的均值为 $\bar{x}_{gj}$，总体均值为 $\bar{x}_{\cdot j}$，组间平方和与组内平方和为
\begin{equation}
SSB_j=\sum_{g\in\mathcal{G}}n_g(\bar{x}_{gj}-\bar{x}_{\cdot j})^2,\qquad
SSW_j=\sum_{g\in\mathcal{G}}\sum_{i:y_i=g}(x_{ij}-\bar{x}_{gj})^2,
\end{equation}
位置 $j$ 的 F 统计量与效应量为
\begin{equation}
F_j=\frac{SSB_j/(G-1)}{SSW_j/(N-G)},\qquad
\eta_j^2=\frac{SSB_j}{SSB_j+SSW_j}.
\end{equation}
可得$F_j$ 为组间均方与组内均方之比，越大说明该位置上类别带来的差异越压过类内噪声，$\eta_j^2\in[0,1]$ 将这一差异折算为比例，表示该位置的总变异中有多大份额由类别造成。两个指标一个看显著性, 一个看效应大小，互为补充。按 $F_j$ 排序的前 10 个检测位置列于表~\ref{tab:anova}，全部前 30 个位置见附录~B。

\begin{table}[H]
\centering
\caption{区分力最强的前 10 个检测位置}
\label{tab:anova}
\begin{tabular}{ccccccc}
\toprule
检测位置 & A 均值 & B 均值 & C 均值 & 类均值极差 & ANOVA $F$ & $\eta^2$ \\
\midrule
134 & 7.0906 & 6.2049 & 3.8486 & 3.2419 & 4731.81 & 0.5427 \\
138 & 7.0406 & 6.2220 & 3.8489 & 3.1917 & 4599.51 & 0.5356 \\
88  & 7.0275 & 6.2341 & 3.8305 & 3.1971 & 4580.39 & 0.5346 \\
108 & 7.0803 & 6.2157 & 3.8488 & 3.2315 & 4534.41 & 0.5321 \\
75  & 7.0778 & 6.1807 & 3.8282 & 3.2495 & 4523.90 & 0.5315 \\
125 & 7.0422 & 6.2153 & 3.8350 & 3.2071 & 4517.32 & 0.5311 \\
78  & 7.0341 & 6.2179 & 3.8323 & 3.2017 & 4516.79 & 0.5311 \\
62  & 7.0495 & 6.2263 & 3.8362 & 3.2133 & 4515.52 & 0.5310 \\
139 & 7.0294 & 6.1828 & 3.8417 & 3.1877 & 4511.47 & 0.5308 \\
89  & 7.0034 & 6.2264 & 3.8319 & 3.1716 & 4467.36 & 0.5284 \\
\bottomrule
\end{tabular}
\end{table}

可知，结果呈现出两个特征。排名前列的位置效应量均在 0.52 以上，一半以上的变异由类别造成，区分力相当强。同时，这些高区分力位置在 155 个位置中分布零散（如 134, 138, 88, 108, 75 并不相邻），且类均值次序一律是 A 高于 B 高于 C，与全局统计特征完全一致。这说明三类差异并非由某几个孤立位置驱动，而是整条检测曲线的系统性平移，几乎每个位置都在以相近的方式将三类拉开。这一发现有两层含义：其一，任意一段曲线都携带类别信息，缺失个别位置不会令分类失效；其二，155 维之间高度冗余，降维具备数据基础。图~\ref{fig:feat} 以柱状形式展示效应量排名。

\begin{figure}[H]
\centering
\includegraphics[width=1.00\textwidth]{04_feature_importance}
\caption{关键检测位置的 $\eta^2$ 排名}
\label{fig:feat}
\end{figure}

\subsection{主成分分析与低维结构}

方差分析提示 155 个位置高度冗余，主成分分析可以将这种冗余加以量化。其在标准化空间中寻找一组互相正交的方向，使数据沿第一个方向的方差最大, 沿第二个方向在与第一个正交的前提下方差次大，依此类推。若前几个方向就吸收了绝大部分方差，就说明高维曲线实际上生活在一个低维子空间里。

对标准化矩阵 $\mathbf{Z}$ 做主成分分析。由于先前标准化操作已使每列均值为 0 且方差为 1，标准化数据的样本协方差矩阵等于检测位置之间的相关矩阵
\begin{equation}
\mathbf{S}=\frac{1}{N-1}\mathbf{Z}^{T}\mathbf{Z}.
\end{equation}
对 $\mathbf{S}$ 做特征分解 $\mathbf{S}\mathbf{v}_k=\lambda_k\mathbf{v}_k$，第 $k$ 个主成分的解释率为 $R_k=\lambda_k/\sum_{\ell}\lambda_\ell$。前三个主成分的解释率分别为 0.5417, 0.0050, 0.0049，前十个累计解释 0.5848。第一主成分单独解释了 54.17\% 的总变异，其余主成分解释率迅速下降并趋于平稳，碎石图见图~\ref{fig:scree}。

\begin{figure}[H]
\centering
\includegraphics[width=1.0\textwidth]{03_pca_scree}
\caption{主成分碎石图与累计解释率}
\label{fig:scree}
\end{figure}

第一主成分独占 54.17\% 的方差，而第二主成分骤降到 0.50\%，这种悬殊说明 155 维曲线的主要变化几乎只沿一个方向发生。结合各位置类均值次序一致的发现，该方向可以理解为整条曲线的综合透射强度，它将每个位置上 A 高 B 中 C 低的同向差异汇聚为一个总指标。其余主成分解释率都在 0.5\% 上下且缓慢衰减，对应的是与类别关系不大的零散噪声。

三类样本在前三个主成分上的类中心进一步印证了这一点。第一主成分上 A 类为 $-7.11$, B 类为 $-1.13$, C 类为 16.10，而第二和第三主成分上的类中心都贴近 0。区分三类的几乎全部信息都压在第一主成分上，后续主成分基本不携带类别信息。图~\ref{fig:scatter} 给出样本在前两个主成分平面上的分布。

\begin{figure}[H]
\centering
\includegraphics[width=1.0\textwidth]{02_pca_scatter}
\caption{前两个主成分平面上的三类工件分布与类中心}
\label{fig:scatter}
\end{figure}

C 类点群独立位于右侧，与 A, B 两类隔开一段明显的空白。A 与 B 两类点群在第一主成分方向上部分交叠，沿第二主成分则几乎完全混合。沿综合强度方向，C 类及容易被区分，A 与 B 却相互混杂，仅靠第一主成分这一个水平方向不足以分开 A 与 B，要区分二者必须回到前面发现的波动差异上。

\section{分类模型建立与求解}
\subsection{判别框架}

本文的分类模型统一建立在贝叶斯判别框架上。给定一件工件的标准化检测向量 $\mathbf{z}$，最优判别规则是将它归入后验概率最大的类别
\begin{equation}
\hat{y}=\arg\max_{g\in\mathcal{G}}P(y=g\mid\mathbf{z}).
\end{equation}
由贝叶斯公式 $P(y=g\mid\mathbf{z})\propto P(\mathbf{z}\mid y=g)\,\pi_g$，后验由类条件似然 $P(\mathbf{z}\mid y=g)$ 与先验 $\pi_g$ 共同决定。对似然取对数不改变其大小次序，因此各模型的判别函数都可写成对数似然加对数先验的形式，差别在于对类条件分布 $P(\mathbf{z}\mid y=g)$ 施加了什么假设。因此后续以强弱不同的分布假设逼近真实的类条件分布，再借助交叉验证比较哪一种假设与数据最为契合。

沿此思路，本文选取五个假设强度递增的模型。最近质心只用类中心，等价于假设各类协方差相同且为单位阵。朴素贝叶斯假设各检测位置在给定类别后相互独立，从而利用位置级方差；线性判别分析假设三类服从协方差相同的多元正态，给出线性判别边界。二次判别分析放开协方差相同这一限制，允许各类拥有各自的协方差矩阵，给出二次判别边界；PCA 降维后的二次判别则在低维空间重做判别，用于检验高维结论是否依赖噪声方向。

\subsection{最近质心模型}

最近质心模型只使用类中心信息。类别 $g$ 的中心为 $\boldsymbol{\mu}_g=\frac{1}{n_g}\sum_{i:y_i=g}\mathbf{z}_i$，样本到各类中心的平方距离为 $d_g^2(\mathbf{z})=\|\mathbf{z}-\boldsymbol{\mu}_g\|_2^2$，判别规则取距离最小的类别
\begin{equation}
\hat{y}_{NC}=\arg\min_{g\in\mathcal{G}}d_g^2(\mathbf{z}).
\end{equation}
从几何上看，最近质心将空间按到三个类中心的远近划分，两类之间的边界是连接两类中心线段的垂直平分面，分界完全由类中心位置决定，与各类的散布形态无关。该模型假设最简, 计算最省，适合用作判断三类是否存在明显类中心分离的基线。其局限在于当两类均值接近而散布形状不同时，垂直平分面会从两团数据中间穿过，将本可凭波动区分的样本判错。A 与 B 恰好符合这种情形，因此最近质心在 A, B 之间的表现可以作为仅用水平信息所能达到的性能参照。

\subsection{高斯朴素贝叶斯}

朴素贝叶斯不再假设所有位置同等散布，而是允许每个位置在每一类里有各自的方差，代价是假设这些位置在给定类别后相互独立。类条件似然因而分解为各位置一维正态密度的乘积，$z_j\mid y=g\sim N(\mu_{gj},\sigma_{gj}^2)$。将似然取对数并去掉与类别无关的常数，加上对数先验，得到判别函数
\begin{equation}
\delta_g^{NB}(\mathbf{z})=\log\pi_g-\frac{1}{2}\sum_{j=1}^{p}\log\sigma_{gj}^2
-\frac{1}{2}\sum_{j=1}^{p}\frac{(z_j-\mu_{gj})^2}{\sigma_{gj}^2}.
\end{equation}
当155 个一维密度直接连乘会迅速下溢为零，化成求和后数值稳定。式中第二项是各类的方差惩罚，第三项是按各位置方差加权的偏差平方和。按方差加权的结构，朴素贝叶斯能够利用 A 类波动大, B 类波动小的差异，而这是最近质心所不具备的。先验取经验比例 $\pi_g=n_g/N$，代入本题为 $\pi_A=0.4477$, $\pi_B=0.3313$, $\pi_C=0.2210$。

相邻检测位置之间往往存在相关性，将它们视为独立会高估似然中的信息量。但只要各位置的边缘分布本身已经能将类别拉开，这种简化带来的偏差就有限，模型仍可表现良好\upcite{handyu2001}。在本题中，该假设的代价并不均匀，后文相关性诊断显示 B 类相邻位置的相关明显高于 A 和 C 两类，独立假设在 B 类上偏离最大，朴素贝叶斯最可能在 B 类上出现误判。因此，朴素贝叶斯的表现可以读作位置级边缘分布能将类别分到何种程度的一个标尺。

\subsection{正则化线性判别分析}

线性判别分析\upcite{fisher1936}将各位置独立这一过强的假设放开，转而假设每一类服从多元正态，且三类共用同一个协方差矩阵：$\mathbf{z}\mid y=g\sim N(\boldsymbol{\mu}_g,\boldsymbol{\Sigma})$。位置之间的相关性由此被纳入模型，但三类的散布形状被强制设为相同。共享协方差由各类离差合并估计
\begin{equation}
\mathbf{S}_p=\frac{1}{N-G}\sum_{g\in\mathcal{G}}\sum_{i:y_i=g}(\mathbf{z}_i-\boldsymbol{\mu}_g)(\mathbf{z}_i-\boldsymbol{\mu}_g)^{T}.
\end{equation}
将多元正态密度代入对数后验，二次项 $\mathbf{z}^T\boldsymbol{\Sigma}^{-1}\mathbf{z}$ 在三类中相同因而相消，判别函数化为关于 $\mathbf{z}$ 的线性函数，类与类之间的边界是一张超平面。

此处有一个高维场景下的现实问题，协方差矩阵 $\mathbf{S}_p$ 有 $155\times156/2\approx1.2$ 万个待估元素，且前面已看到相邻位置高度相关，直接求逆会因接近奇异而极不稳定，估计出的判别方向可能被噪声放大。常用的补救手段是将样本协方差向其对角阵收缩\upcite{friedman1989}，即在完全保留样本相关与只保留各位置方差之间取折中
\begin{equation}
\boldsymbol{\Sigma}_{LDA}=(1-\lambda)\mathbf{S}_p+\lambda\operatorname{diag}(\mathbf{S}_p)+\epsilon\mathbf{I},
\end{equation}
其中 $\epsilon=10^{-6}$ 仅用于保证正定，收缩参数 $\lambda$ 越大越偏向对角, 越稳健但也越接近朴素贝叶斯。$\lambda$ 在候选集合 $\{0,0.02,0.05,0.08,0.12,0.20,0.40,0.60\}$ 上做网格敏感性分析（见表~\ref{tab:shrink}），准确率随 $\lambda$ 增大略有上升并在 $\lambda=0.60$ 处取得最高，故取 $\lambda=0.60$。判别函数为
\begin{equation}
\delta_g^{LDA}(\mathbf{z})=\mathbf{z}^{T}\boldsymbol{\Sigma}_{LDA}^{-1}\boldsymbol{\mu}_g
-\frac{1}{2}\boldsymbol{\mu}_g^{T}\boldsymbol{\Sigma}_{LDA}^{-1}\boldsymbol{\mu}_g+\log\pi_g.
\end{equation}
由于边界是线性的，LDA 能否分开 A 与 B，直接回答了一张超平面是否足够的问题。考虑到 A 与 B 的差异主要在波动而非水平，而线性边界本质上只在比较加权后的水平，可以预期 LDA 在 A 与 B 之间的区分能力有限。本文用它作为线性可分性的基准。

\subsection{正则化二次判别分析}

A 与 B 的关键差异在散布而非位置，线性判别分析受共享协方差假设的约束，恰恰无法捕捉这种散布差异。二次判别分析假设每一类有各自的协方差 $\mathbf{z}\mid y=g\sim N(\boldsymbol{\mu}_g,\boldsymbol{\Sigma}_g)$，于是不同类的二次项不再相消，判别边界变成二次曲面，能够表达某类整体更分散这类形状差异。与 LDA 同理，每类协方差按其样本量估计，仍需对角收缩以保证数值稳定
\begin{equation}
\boldsymbol{\Sigma}_g^{*}=(1-\alpha)\mathbf{S}_g+\alpha\operatorname{diag}(\mathbf{S}_g)+\epsilon\mathbf{I},
\end{equation}
其中 $\mathbf{S}_g$ 为类别 $g$ 的样本协方差。判别函数为
\begin{equation}
\delta_g^{QDA}(\mathbf{z})=-\frac{1}{2}(\mathbf{z}-\boldsymbol{\mu}_g)^{T}(\boldsymbol{\Sigma}_g^{*})^{-1}(\mathbf{z}-\boldsymbol{\mu}_g)
-\frac{1}{2}\log|\boldsymbol{\Sigma}_g^{*}|+\log\pi_g.
\label{eq:qda}
\end{equation}
为保证行列式与逆矩阵的数值稳定，对数行列式通过特征值计算
\begin{equation}
\log|\boldsymbol{\Sigma}_g^{*}|=\sum_{r}\log\max(\lambda_{gr},10^{-12}),
\end{equation}
其中 $\lambda_{gr}$ 为 $\boldsymbol{\Sigma}_g^{*}$ 的第 $r$ 个特征值。判别函数~\eqref{eq:qda}由三部分构成，各有明确的统计含义。第一项是样本到类中心的马氏距离，按该类自身的协方差度量样本距该类的远近，第二项 $-\frac{1}{2}\log|\boldsymbol{\Sigma}_g^{*}|$ 是体积惩罚，协方差越大的类其分布越弥散，容易将远处的点也吸纳进来，该项对此施加修正，第三项是对数先验。正是第一项中每类使用各自的协方差度量距离，使 QDA 能够借助 A 类散布更大这一特征将落在 A, B 之间的样本正确归类。

收缩参数 $\alpha$ 在同一网格上做敏感性分析（见表~\ref{tab:shrink}），所有候选取值在交叉验证中均达到满分，说明 QDA 的满分并不依赖特定的收缩取值，本文取 $\alpha=0.12$ 以兼顾协方差矩阵的数值稳定。式~\eqref{eq:qda}同时利用类别均值差异与类别内部协方差结构差异，与特征分析得到的 A, B 差在波动这一结论完全对应，因此本文以 QDA 为主分类器。

\subsection{PCA-QDA 低维复核}

155 维 QDA 在每类需估计一个 $155\times155$ 的协方差矩阵，参数众多，一个合理的顾虑是其是否混杂了与类别无关的噪声方向而过拟合。因此后续使用PCA-QDA检验，即先在训练集内拟合 PCA 取前 $r$ 个主成分构成投影矩阵 $\mathbf{V}_r$，将样本压到低维空间 $\mathbf{u}_i=\mathbf{V}_r^{T}(\mathbf{z}_i-\bar{\mathbf{z}})$，再在该空间重做正则化 QDA（低维收缩取 $\alpha=0.10$）。低维空间丢掉了高维噪声方向，若在此仍能得到与高维一致的判别结果，就说明结论依赖的是数据中的真实结构而非噪声。维数 $r$ 不取某个最优值，而作为稳健性旋钮在 $\{1,2,3,5,8,10,12,20,30,50\}$ 上扫描：只要未知工件的判别在如此宽的维数范围内都不变，分类结论就不取决于人为设定的维数。降维方法的依据见 PCA 的综述文献\upcite{jolliffe2016}。

\subsection{一致性复核与拒判规则}

记五个模型的集合为 $\mathcal{M}=\{NC,NB,LDA,QDA,\text{PCA-QDA}\}$，每个模型对样本给出一个预测 $h_m(\mathbf{z})$。由于主分类器 QDA 在交叉验证中已达到满分，多模型组合在此不提供独立的性能提升，其定位是一致性复核。综合判别取五个模型预测的众数，即得票最多的类别
\begin{equation}
\hat{y}_{ens}=\arg\max_{g\in\mathcal{G}}\sum_{m\in\mathcal{M}}\mathbb{1}[h_m(\mathbf{z})=g].
\end{equation}
模型一致数 $C(\mathbf{z})=\sum_{m}\mathbb{1}[h_m(\mathbf{z})=\hat{y}_{ens}]$ 表示与最终判别一致的模型个数，用于度量分类可信度。当各模型出现分歧时仍以 QDA 为主判依据，因其在交叉验证中取得最高的平衡准确率。

考虑到生产中可能出现 A, B, C 之外的新类型或设备漂移，本文给出拒判规则，即当模型一致数低于 4，或 QDA 的最大后验概率低于阈值 $\tau_P$ 时，将样本标记为待人工复检
\begin{equation}
C(\mathbf{z})<4\quad\text{或}\quad P_{\max}^{QDA}(\mathbf{z})<\tau_P,
\end{equation}
其中 $\tau_P$ 可初始取 0.80 并用历史复检样本校准。该规则使模型在面对异常工件时不会强行给出分类。

\subsection{评价指标}

对类别集合 $\mathcal{G}=\{A,B,C\}$，混淆矩阵 $\mathbf{C}=(c_{gh})$ 的元素 $c_{gh}=\sum_{i}\mathbb{1}(y_i=g,\hat{y}_i=h)$ 表示真实类别为 $g$ 却被判为 $h$ 的样本数，矩阵行为真实类别, 列为预测类别，对角元越大说明判别越准确。由混淆矩阵定义总准确率, 类别召回率与类别精确率
\begin{equation}
ACC=\frac{1}{N}\sum_{i=1}^{N}\mathbb{1}(\hat{y}_i=y_i),\qquad
Recall_g=\frac{c_{gg}}{\sum_{h}c_{gh}},\qquad
Precision_g=\frac{c_{gg}}{\sum_{h}c_{hg}}.
\end{equation}
召回率衡量某类样本被正确识别的比例，精确率衡量判为某类的样本中真正属于该类的比例。两者的调和平均为类别 F1，对各类取算术平均得到宏 F1 与宏精确率
\begin{equation}
F1_g=\frac{2\,Precision_g\,Recall_g}{Precision_g+Recall_g},\qquad
\text{Macro-}F1=\frac{1}{G}\sum_{g\in\mathcal{G}}F1_g.
\end{equation}
由于三类样本数量不等，再引入对各类等权的平衡准确率
\begin{equation}
BA=\frac{1}{G}\sum_{g\in\mathcal{G}}Recall_g,
\end{equation}
它避免样本量最大的 A 类主导评价，能更公平地反映样本较少的 C 类的识别情况。

\subsection{求解流程}

对 10 件未知工件，先用训练集的均值与标准差将它们映射到同一标准化空间，再分别由五个模型给出预测。对每件工件，取五个预测的众数作为最终类别，并记录与最终类别一致的模型个数。当一致数低于 4 或 QDA 的最大后验概率低于阈值时，不直接输出类别而转人工复检。整体建模流程见图~\ref{fig:flow}，判别结果在结果分析部分给出。

\section{结果分析}
\subsection{特征分析结果}

特征分析回答了题目的第一项要求。我们可以发现三类工件的差异由两部分构成。整体透射水平的差异使 C 类与 A, B 两类明显分离，曲线波动程度的差异使 A 类区别于 B 和 C 两类。方差分析定位出一批效应量超过 0.52 的关键检测位置，主成分分析进一步表明类别差异主要沿第一主成分这一综合强度方向展开。这些结果共同说明三类工件在 155 维检测空间中具有清晰的可分结构，为后续分类提供了实证依据。

\subsection{模型性能比较}

五个模型在分层五折交叉验证下的性能列于表~\ref{tab:cv}，对应的总准确率, 平衡准确率与宏 F1 见图~\ref{fig:perf}。

\begin{table}[H]
\centering
\caption{五个统计判别模型的交叉验证性能}
\label{tab:cv}
\resizebox{\textwidth}{!}{%
\begin{tabular}{lccccccc}
\toprule
模型 & 总准确率 & 宏精确率 & 平衡准确率 & 宏 F1 & A 召回 & B 召回 & C 召回 \\
\midrule
最近质心        & 0.9535 & 0.9626 & 0.9564 & 0.9589 & 0.9731 & 0.8960 & 1.0000 \\
高斯朴素贝叶斯  & 0.9917 & 0.9940 & 0.9917 & 0.9927 & 1.0000 & 0.9750 & 1.0000 \\
正则化 LDA      & 0.9506 & 0.9609 & 0.9533 & 0.9563 & 0.9748 & 0.8850 & 1.0000 \\
正则化 QDA      & 1.0000 & 1.0000 & 1.0000 & 1.0000 & 1.0000 & 1.0000 & 1.0000 \\
PCA-QDA(12 维)  & 0.9580 & 0.9628 & 0.9639 & 0.9633 & 0.9471 & 0.9448 & 1.0000 \\
\bottomrule
\end{tabular}}
\end{table}

\begin{figure}[H]
\centering
\includegraphics[width=0.92\textwidth]{06_model_performance}
\caption{五个统计判别模型的性能对比}
\label{fig:perf}
\end{figure}
最近质心与正则化 LDA 都只利用了水平信息，总准确率停在 0.95 附近，朴素贝叶斯加入各位置的方差后，总准确率跃升至 0.9917，正则化 QDA 进一步利用各类各自的协方差，达到 1.0000。我们可以发现A 与 B 的差异藏在方差结构中，模型对方差信息利用得越充分，分类性能越高。朴素贝叶斯虽然假设各位置独立，却已逼近满分，这说明位置级的方差差异本身就携带了绝大部分类别信息，位置间的相关性在此之上提供的增量有限。

三类的分类别精确率, 召回率与 F1 列于表~\ref{tab:perclass}，将总体性能拆解到每一类上可以观察得更为细致。C 类召回率在所有模型中都是 1.0000，与其在主成分图上独处一隅的位置一致，任何模型都不会漏判 C 类。性能差距在A 与 B的判别之间较为明显，最近质心与 LDA 的 B 类召回率只有 0.8960 与 0.8850，约一成的真 B 被误判，朴素贝叶斯将 B 类召回率提升到 0.9750，QDA 则达到 1.0000。PCA-QDA 在 12 维时总准确率回落到 0.9580，原因在于只保留 12 维时丢掉了部分用于区分 A, B 的方差信息，用于后文的稳健性复核。

\begin{table}[H]
\centering
\caption{各模型的分类别精确率、召回率与 F1}
\label{tab:perclass}
\resizebox{\textwidth}{!}{
\begin{tabular}{lccccccccc}
\toprule
\multirow{2}{*}{模型} & \multicolumn{3}{c}{A 类} & \multicolumn{3}{c}{B 类} & \multicolumn{3}{c}{C 类} \\
\cmidrule(lr){2-4}\cmidrule(lr){5-7}\cmidrule(lr){8-10}
 & 精确率 & 召回率 & F1 & 精确率 & 召回率 & F1 & 精确率 & 召回率 & F1 \\
\midrule
最近质心       & 0.9267 & 0.9731 & 0.9493 & 0.9610 & 0.8960 & 0.9274 & 1.0000 & 1.0000 & 1.0000 \\
高斯朴素贝叶斯 & 0.9819 & 1.0000 & 0.9908 & 1.0000 & 0.9750 & 0.9874 & 1.0000 & 1.0000 & 1.0000 \\
正则化 LDA     & 0.9197 & 0.9748 & 0.9465 & 0.9629 & 0.8850 & 0.9223 & 1.0000 & 1.0000 & 1.0000 \\
正则化 QDA     & 1.0000 & 1.0000 & 1.0000 & 1.0000 & 1.0000 & 1.0000 & 1.0000 & 1.0000 & 1.0000 \\
PCA-QDA(12 维) & 0.9586 & 0.9471 & 0.9528 & 0.9296 & 0.9448 & 0.9371 & 1.0000 & 1.0000 & 1.0000 \\
\bottomrule
\end{tabular}}
\end{table}

\subsection{混淆结构与 A, B 方差可分机制}

图~\ref{fig:cm} 给出朴素贝叶斯, LDA 与 QDA 的混淆矩阵。然而三个模型有一个共同点，即误判全部落在 A 与 B 之间，没有任何样本与 C 类发生混淆，再次印证了 C 类的独立性。误判还呈现明显的方向性，主要表现为真 B 被判成 A，LDA 将 304 件真 B 判为 A，朴素贝叶斯将 66 件真 B 判为 A。方向性可以这样理解：A 的散布比 B 大，按马氏距离衡量时 A 类的覆盖范围更宽，靠近边界的 B 样本更容易落入 A 的覆盖区域，只使用水平信息或线性边界的模型尤其如此。QDA 由于分别度量到两类的马氏距离并扣除体积惩罚，纠正了这种偏向，三类上均无误判。

\begin{figure}[H]
\centering
\includegraphics[width=0.98\textwidth]{07_confusion_matrices}
\caption{朴素贝叶斯、LDA 与 QDA 的混淆矩阵}
\label{fig:cm}
\end{figure}

QDA 取得满分而 LDA 与最近质心不能，根本原因在于 A 与 B 的差异主要体现在方差而非均值。为量化这一点，对 A, B 两类的若干一维特征计算分布重叠系数与最优单阈值准确率，结果列于表~\ref{tab:overlap}。

\begin{table}[H]
\centering
\caption{A, B 两类在不同特征上的重叠与单阈值可分性}
\label{tab:overlap}
\begin{tabular}{lcc}
\toprule
特征 & 直方图重叠系数 & 最优单阈值准确率 \\
\midrule
样本均值        & 0.1139 & 0.9426 \\
样本内标准差    & 0.0000 & 1.0000 \\
样本极差        & 0.0095 & 0.9955 \\
样本四分位距    & 0.0066 & 0.9969 \\
第一主成分得分  & 0.1143 & 0.9426 \\
前 15 关键位置均值 & 0.3025 & 0.8484 \\
\bottomrule
\end{tabular}
\end{table}

结果分成对比鲜明的两组。刻画水平的特征重叠系数约 0.11，单阈值准确率都只有约 0.9426，与最近质心和 LDA 的总体性能吻合，说明这些模型触及的正是同一个性能天花板：水平信息无法分开 A 与 B。刻画波动的特征则呈现截然不同的面貌，样本内标准差的重叠系数为 0, 极差与四分位距分别低至 0.0095 与 0.0066，单阈值准确率达到 0.9955 以上。其中样本内标准差的分离最为彻底，重叠为 0, 单阈值准确率为 1.0000。得其最优切割点为
\begin{equation}
s^{*}=1.105855,
\end{equation}
判别规则为 $s_i>s^{*}$ 判为 A 类，否则判为 B 类。训练数据中 A 类样本内标准差的最小值为 1.139050，B 类的最大值为 1.072661，两者被阈值完全隔开，该一维规则在训练集上实现了 A 与 B 的完全分离。图~\ref{fig:abstd} 有直观展示。

\begin{figure}[H]
\centering
\includegraphics[width=0.88\textwidth]{05_ab_std_threshold}
\caption{A, B 两类样本内标准差分布与最优阈值}
\label{fig:abstd}
\end{figure}

QDA 通过类别协方差结构捕捉到 A 类波动显著大于 B 类的特征，而 LDA 假设两类协方差相同，无法利用这一差异。该结论也提示，QDA 的优势依赖检测设备的方差标定保持稳定，一旦扫描噪声水平或尺度标定发生变化，A, B 的方差差异可能被压缩，模型需要重新校准。

\subsection{未知工件分类结果}

10 件未知工件的判别结果列于表~\ref{tab:unknown}。五个模型对每件工件给出完全一致的判别，模型一致数均为 5。表中 QDA 最大后验概率由判别分数经 softmax 归一化得到并四舍五入为四位小数，10 件样本均为 1.0000，表明各样本到所属类别的判别分数显著高于其他两类。

\begin{table}[H]
\centering
\caption{未知工件分类结果}
\label{tab:unknown}
\resizebox{\textwidth}{!}{%
\begin{tabular}{ccccccccc}
\toprule
工件 & 最终类别 & 一致数 & 最近质心 & 朴素贝叶斯 & LDA & QDA & PCA-QDA & QDA 最大后验概率 \\
\midrule
D1  & A & 5/5 & A & A & A & A & A & 1.0000 \\
D2  & A & 5/5 & A & A & A & A & A & 1.0000 \\
D3  & A & 5/5 & A & A & A & A & A & 1.0000 \\
D4  & B & 5/5 & B & B & B & B & B & 1.0000 \\
D5  & B & 5/5 & B & B & B & B & B & 1.0000 \\
D6  & B & 5/5 & B & B & B & B & B & 1.0000 \\
D7  & C & 5/5 & C & C & C & C & C & 1.0000 \\
D8  & C & 5/5 & C & C & C & C & C & 1.0000 \\
D9  & C & 5/5 & C & C & C & C & C & 1.0000 \\
D10 & A & 5/5 & A & A & A & A & A & 1.0000 \\
\bottomrule
\end{tabular}}
\end{table}

最终判别为：D1, D2, D3, D10 属于 A 类，D4, D5, D6 属于 B 类，D7, D8, D9 属于 C 类。五个假设各异的模型给出一致结论，且 QDA 判别分数差距显著，说明在当前数据分布与交叉验证协议下，这 10 件工件均不处于类别边界附近，判别结果的可信度较高。

\section{模型检验与稳健性分析}
\subsection{交叉验证与各折稳定性}

为防止仅报告总体准确率而使评价过度乐观，表~\ref{tab:fold} 列出五折交叉验证的各折准确率。QDA 在五折上均为 1.0000，并非某一次划分的偶然结果。其余模型各折准确率的折间极差均在约 1 个百分点以内。评价采用分层抽样保证每折的类别比例与总体一致\upcite{kohavi1995}。

\begin{table}[H]
\centering
\caption{五折交叉验证的各折准确率}
\label{tab:fold}
\begin{tabular}{lccccc}
\toprule
模型 & 第 1 折 & 第 2 折 & 第 3 折 & 第 4 折 & 第 5 折 \\
\midrule
最近质心       & 0.9549 & 0.9537 & 0.9486 & 0.9542 & 0.9561 \\
高斯朴素贝叶斯 & 0.9894 & 0.9906 & 0.9906 & 0.9944 & 0.9937 \\
正则化 LDA     & 0.9537 & 0.9499 & 0.9436 & 0.9529 & 0.9529 \\
正则化 QDA     & 1.0000 & 1.0000 & 1.0000 & 1.0000 & 1.0000 \\
PCA-QDA(12 维) & 0.9580 & 0.9574 & 0.9536 & 0.9624 & 0.9586 \\
\bottomrule
\end{tabular}
\end{table}

\subsection{先验概率敏感性}

主模型采用经验先验。为检验先验选择对结论的影响，将经验先验与等先验对比，结果列于表~\ref{tab:prior}。先验变化对朴素贝叶斯与 LDA 仅产生微小影响，对 QDA 没有任何影响，未知工件的判别结果亦不改变。若训练比例不能代表生产比例，可在部署时改用等先验或生产历史比例先验。

\begin{table}[H]
\centering
\caption{先验概率敏感性}
\label{tab:prior}
\begin{tabular}{llcc}
\toprule
模型 & 先验 & 总准确率 & 平衡准确率 \\
\midrule
朴素贝叶斯 & 等先验   & 0.9919 & 0.9918 \\
朴素贝叶斯 & 经验先验 & 0.9917 & 0.9917 \\
LDA       & 等先验   & 0.9521 & 0.9552 \\
LDA       & 经验先验 & 0.9506 & 0.9533 \\
QDA       & 等先验   & 1.0000 & 1.0000 \\
QDA       & 经验先验 & 1.0000 & 1.0000 \\
\bottomrule
\end{tabular}
\end{table}

\subsection{PCA 维数敏感性}

PCA-QDA 在不同维数下的性能列于表~\ref{tab:pcadim}。$r=12$ 并非最优维数，取该值仅作为低维稳健性检查点。提高维数能小幅改善验证准确率，但仍未达到原始 155 维 QDA 的水平。更重要的是，在全部考察的维数下，10 件未知工件的判别结果与主 QDA 完全一致，说明分类结论不依赖人为设定的降维维数。

\begin{table}[H]
\centering
\caption{PCA-QDA 在不同主成分维数下的性能}
\label{tab:pcadim}
\resizebox{\textwidth}{!}{%
\begin{tabular}{ccccccccccc}
\toprule
维数 $r$ & 1 & 2 & 3 & 5 & 8 & 10 & 12 & 20 & 30 & 50 \\
\midrule
总准确率   & 0.9542 & 0.9553 & 0.9565 & 0.9566 & 0.9574 & 0.9570 & 0.9580 & 0.9554 & 0.9581 & 0.9799 \\
平衡准确率 & 0.9589 & 0.9601 & 0.9615 & 0.9619 & 0.9629 & 0.9628 & 0.9639 & 0.9624 & 0.9647 & 0.9831 \\
\bottomrule
\end{tabular}}
\end{table}

\subsection{参数稳健性}

对两个含收缩参数的判别模型做网格敏感性分析，各收缩取值下的分层五折交叉验证性能列于表~\ref{tab:shrink}。LDA 的总准确率随 $\lambda$ 增大的变化幅度在小数点后第三位，$\lambda=0.60$ 处取得最高的 0.9506，故选此值。QDA 在整个网格上的总准确率与平衡准确率均为 1.0000，说明其满分不依赖特定收缩取值，取 $\alpha=0.12$ 是出于协方差矩阵的数值稳定考虑。需要说明的是，标准化参数与 PCA 投影都在每个训练折内部估计后再作用于验证折，从而避免验证集信息通过这些步骤泄漏到训练过程\upcite{cawley2010}。由于三类样本数量不等，评价同时报告平衡准确率，使样本量最少的 C 类与样本量最多的 A 类获得同等权重\upcite{brodersen2010}。

\begin{table}[H]
\centering
\caption{LDA 与 QDA 的收缩参数网格敏感性}
\label{tab:shrink}
\begin{tabular}{ccccc}
\toprule
收缩参数 & LDA 总准确率 & LDA 平衡准确率 & QDA 总准确率 & QDA 平衡准确率 \\
\midrule
0.00 & 0.9477 & 0.9503 & 1.0000 & 1.0000 \\
0.02 & 0.9480 & 0.9506 & 1.0000 & 1.0000 \\
0.05 & 0.9481 & 0.9507 & 1.0000 & 1.0000 \\
0.08 & 0.9486 & 0.9511 & 1.0000 & 1.0000 \\
0.12 & 0.9489 & 0.9515 & 1.0000 & 1.0000 \\
0.20 & 0.9486 & 0.9513 & 1.0000 & 1.0000 \\
0.40 & 0.9502 & 0.9529 & 1.0000 & 1.0000 \\
0.60 & 0.9506 & 0.9533 & 1.0000 & 1.0000 \\
\bottomrule
\end{tabular}
\end{table}

\section{模型评价与改进}
\subsection{模型优点}

本文方法最核心的优点在于可解释性，从特征分析到判别建模再到部署监控，每一步都对应一个含义明确的统计量。方差分析指出哪些检测位置区分力强，主成分分析表明差异沿哪个方向展开，二次判别分析的马氏距离与体积惩罚解释了 A, B 能够分开的原因。判别边界均有显式公式，结论可以逐项复核。

在验证层面，本文的评价覆盖了分层交叉验证, 各折准确率, 混淆矩阵, 分类别指标与平衡准确率，并辅以先验, 收缩参数与降维维数三类敏感性分析，使结论的来源与适用边界都有据可查。

在稳健性层面，10 件未知工件在五个假设各异的模型下判别完全相同，主 QDA 与 $r$ 从 1 到 50 的 PCA-QDA 也给出相同判别，说明结论既不依赖某个模型的特定偏好，也不取决于某个维数的人为设定。

\subsection{模型局限}

本文的验证建立在单一数据集上，没有跨批次或跨设备的外部测试，上述结论的可迁移性仍需在真实产线上进一步检验。此外还有两点结构性局限。第一，训练数据只包含 A, B, C 三类，若未来出现新型号工件，模型仍会将其强行归入三类之一，必须依靠拒判与人工复检兜底。第二，本文将 155 个检测位置作为无序向量处理，没有利用位置之间的空间关系；按类别计算时 A, B, C 三类相邻位置的平均相关系数分别约为 0.1189, 0.6679, 0.0884，B 类内部仍保留较强的相邻相关，若已知检测位置的几何排列，引入曲线平滑或函数型数据分析有望进一步压缩噪声, 提升跨批次稳健性。

\subsection{生产部署的拒判与漂移监控}

为支撑模型在真实产线的长期运行，本文给出两项部署建议。一是前述的拒判规则，当模型一致数低于 4 或 QDA 后验概率低于阈值时输出待复检标记。二是批次漂移监控。直接在原始空间用欧氏距离监控会被量值大的位置主导，因此应先用训练集参数将新批次标准化，再在标准化空间用正则化 Hotelling 统计量\upcite{jiangfeng2022}
\begin{equation}
T^2_{batch}=n_b(\bar{\mathbf{z}}_{batch}-\bar{\mathbf{z}}_{train})^{T}(\mathbf{S}_z+\eta\mathbf{I})^{-1}(\bar{\mathbf{z}}_{batch}-\bar{\mathbf{z}}_{train}),
\end{equation}
其中 $n_b$ 为批次样本数，$\bar{\mathbf{z}}_{train}=\mathbf{0}$，$\eta$ 为小的正则化常数。控制限可由历史稳定生产期的批次 $T^2$ 取 99\% 分位数确定。当批次样本量较小时，理论近似不稳定，宜优先采用经验控制限。考虑到第一主成分解释了 54.17\% 的变异，还可同步监控批次的第一主成分均值偏移 $D_{PC1}=|\bar{t}_{batch,1}-\bar{t}_{train,1}|$。当 $T^2_{batch}$ 或 $D_{PC1}$ 异常时，提示扫描尺度或工件分布可能发生漂移，应重新估计标准化参数与协方差矩阵。

\section{结论}

本文针对三类工件 X 光检测数据的特征分析与分类问题，建立了以正则化二次判别分析为主体, 以多种统计判别模型为复核的分类框架。特征分析表明，三类工件的整体透射水平呈 A, B, C 递减，C 类与 A, B 两类明显分离，而 A 与 B 的关键差异在于曲线波动程度。正则化 QDA 在分层五折交叉验证中的各项指标均为 1.0000，其判别力来自类别方差与协方差结构的差异，机制分析给出了 A 与 B 在样本内标准差上完全可分的阈值证据。对 10 件未知工件，五个模型给出一致判别：D1, D2, D3, D10 为 A 类，D4, D5, D6 为 B 类，D7, D8, D9 为 C 类。敏感性分析显示该结论对先验, 正则化参数与降维维数均不敏感。本文进一步给出拒判规则与 Hotelling $T^2$ 漂移监控方案，为模型在生产环境中的部署提供了风险预警机制。

%================== 参考文献 ==================
\clearpage
\begin{thebibliography}{99}
\addcontentsline{toc}{section}{参考文献}
\bibitem{mery2015} Mery D，Computer Vision for X-Ray Testing: Imaging, Systems, Image Databases, and Algorithms，Cham：Springer，2015。

\bibitem{zhanghui2022} 张辉，张邹铨，陈煜嵘，吴天月，钟杭，王耀南，工业铸件缺陷无损检测技术的应用进展与展望，自动化学报，48(4)：935--956，2022。

\bibitem{fisher1936} Fisher R A，The use of multiple measurements in taxonomic problems，Annals of Eugenics，7(2)：179--188，1936。

\bibitem{friedman1989} Friedman J H，Regularized discriminant analysis，Journal of the American Statistical Association，84(405)：165--175，1989。

\bibitem{handyu2001} Hand D J，Yu K，Idiot's Bayes—Not So Stupid After All?，International Statistical Review，69(3)：385--398，2001。

\bibitem{jolliffe2016} Jolliffe I T，Cadima J，Principal component analysis: a review and recent developments，Philosophical Transactions of the Royal Society A，374：20150202，2016。

\bibitem{kohavi1995} Kohavi R，A study of cross-validation and bootstrap for accuracy estimation and model selection，Proceedings of the 14th International Joint Conference on Artificial Intelligence (IJCAI)，2：1137--1143，1995。

\bibitem{cawley2010} Cawley G C，Talbot N L C，On over-fitting in model selection and subsequent selection bias in performance evaluation，Journal of Machine Learning Research，11：2079--2107，2010。

\bibitem{brodersen2010} Brodersen K H，Ong C S，Stephan K E，Buhmann J M，The balanced accuracy and its posterior distribution，Proceedings of the 20th International Conference on Pattern Recognition (ICPR)，3121--3124，2010。

\bibitem{jiangfeng2022} 姜云卢，丰之韵，高维稳健 Hotelling $T^2$ 控制图的研究与应用，系统科学与数学，42(7)：1877--1890，2022。
\end{thebibliography}

%================== 附录 ==================
\newpage
\section*{\appentry{A}{完整建模代码}}
\addcontentsline{toc}{section}{附录 A 建模代码}

\lstset{language=Python,keywordstyle=\color{blue}\bfseries,commentstyle=\color{comment},stringstyle=\color{string},showstringspaces=false,breaklines=true,extendedchars=true,basicstyle=\ttfamily\scriptsize}
\lstinputlisting[language=Python]{../process/final_stepwise_code.py}

\newpage
\section*{\appentry{B}{补充结果表}}
\addcontentsline{toc}{section}{附录 B 补充结果表}
\setcounter{table}{0}
\renewcommand{\thetable}{B-\arabic{table}}

表~\ref{tab:anova30} 给出方差分析中区分力最强的前 30 个检测位置，作为正文表~\ref{tab:anova} 的补充。

\begin{longtable}{ccccccc}
\caption{区分力最强的前 30 个检测位置}\label{tab:anova30}\\
\toprule
检测位置 & A 均值 & B 均值 & C 均值 & 类均值极差 & ANOVA $F$ & $\eta^2$ \\
\midrule
\endfirsthead
\multicolumn{7}{c}{{\tablename\ \thetable{} 续页}}\\
\toprule
检测位置 & A 均值 & B 均值 & C 均值 & 类均值极差 & ANOVA $F$ & $\eta^2$ \\
\midrule
\endhead
\midrule
\multicolumn{7}{r}{{续下页}}\\
\endfoot
\bottomrule
\endlastfoot
134 & 7.0906 & 6.2049 & 3.8486 & 3.2419 & 4731.81 & 0.5427 \\
138 & 7.0406 & 6.2220 & 3.8489 & 3.1917 & 4599.51 & 0.5356 \\
88  & 7.0275 & 6.2341 & 3.8305 & 3.1971 & 4580.39 & 0.5346 \\
108 & 7.0803 & 6.2157 & 3.8488 & 3.2315 & 4534.41 & 0.5321 \\
75  & 7.0778 & 6.1807 & 3.8282 & 3.2495 & 4523.90 & 0.5315 \\
125 & 7.0422 & 6.2153 & 3.8350 & 3.2071 & 4517.32 & 0.5311 \\
78  & 7.0341 & 6.2179 & 3.8323 & 3.2017 & 4516.79 & 0.5311 \\
62  & 7.0495 & 6.2263 & 3.8362 & 3.2133 & 4515.52 & 0.5310 \\
139 & 7.0294 & 6.1828 & 3.8417 & 3.1877 & 4511.47 & 0.5308 \\
89  & 7.0034 & 6.2264 & 3.8319 & 3.1716 & 4467.36 & 0.5284 \\
23  & 7.0582 & 6.2227 & 3.8268 & 3.2314 & 4461.86 & 0.5281 \\
135 & 7.0528 & 6.2149 & 3.8202 & 3.2326 & 4460.78 & 0.5280 \\
19  & 7.0436 & 6.2392 & 3.8245 & 3.2191 & 4456.18 & 0.5278 \\
95  & 7.0343 & 6.2171 & 3.8374 & 3.1969 & 4453.49 & 0.5276 \\
1   & 7.0457 & 6.2219 & 3.8733 & 3.1725 & 4449.76 & 0.5274 \\
109 & 7.0401 & 6.2035 & 3.8712 & 3.1690 & 4447.23 & 0.5273 \\
107 & 7.0396 & 6.2333 & 3.8503 & 3.1893 & 4447.22 & 0.5273 \\
10  & 7.0255 & 6.2124 & 3.8299 & 3.1956 & 4443.93 & 0.5271 \\
145 & 7.0248 & 6.2341 & 3.8448 & 3.1800 & 4442.87 & 0.5270 \\
61  & 7.0278 & 6.2470 & 3.8348 & 3.1930 & 4435.79 & 0.5266 \\
7   & 7.0616 & 6.2284 & 3.8578 & 3.2037 & 4429.15 & 0.5262 \\
63  & 7.0363 & 6.2280 & 3.8743 & 3.1621 & 4415.87 & 0.5255 \\
112 & 7.0564 & 6.2012 & 3.8537 & 3.2027 & 4414.00 & 0.5254 \\
5   & 7.0497 & 6.2201 & 3.8350 & 3.2147 & 4411.85 & 0.5253 \\
73  & 7.0363 & 6.1827 & 3.8647 & 3.1716 & 4410.70 & 0.5252 \\
122 & 7.0603 & 6.2151 & 3.8660 & 3.1944 & 4407.25 & 0.5250 \\
22  & 7.0617 & 6.2372 & 3.8610 & 3.2006 & 4401.27 & 0.5247 \\
47  & 7.0120 & 6.2036 & 3.8526 & 3.1594 & 4400.71 & 0.5246 \\
42  & 7.0804 & 6.2089 & 3.8752 & 3.2051 & 4399.82 & 0.5246 \\
128 & 7.0244 & 6.2408 & 3.8404 & 3.1841 & 4399.65 & 0.5246 \\
\end{longtable}

\end{document}
